% SI numbering: first-citation order (tools/si_numbering.py, round 6, 2026-10-05)
% methods.tex : 방법(국문판). 영문판 paper/manuscript/en/sections/methods.tex(2026-10-05 2차 수정판, XA 계열 수 정정 포함)를
% 문장 단위로 옮겼다. 수식, 수치, 인용 키, 커밋 번호, 자리표시 문자열은 영문판과 같다. 코드 경로는 원문 그대로 둔다.

\section{방법}\label{sec:methods}

\subsection{연구 지역과 라벨}\label{sec:m-labels}

주 지역의 활동층 두께(ALT) 라벨은 GTN-P/CALM 기록\cite{Streletskiy2025}, 레나강 삼각주 지역의 ALLena 자료\cite{Veremeeva2025}, 알래스카와 캐나다의 ABoVE 자료 2판\cite{Moore2025}에서 가져왔다. 한 위치의 반복 측정이 독립 라벨로 들어가지 않도록 기록을 셀로 모았다. 셀은 소수 넷째 자리로 반올림한 좌표(ABoVE와 ALLena)이거나 CALM 지점 하나이며, 셀의 라벨은 그 관측의 평균이다(보충 방법~1). 셀 표는 17,572행이고, 그 가운데 17,467행이 기계식 탐침이나 지표 투과 레이더(ground-penetrating radar, GPR)로 잰 직접 라벨이다. 메타자료에 지중 온도나 융해 관(thaw tube)이 적힌 CALM 셀 68개는 라벨로 쓰지 않는다. 지온에서 유도한 셀 20개(알래스카 15, 캐나다 4, 티베트 고원 1)는 표를 고정한 뒤에 발견되어 표지를 단 채 직접 라벨 범주에 남아 있다. 라벨 연도는 1990--2024년이므로, 예측은 정적이며 2015--2020년 공변량 기후값을 기준으로 한다.

광역 지역은 분석 전에 지리적으로 정하였으며, 알래스카, 레나델타, 캐나다, 러시아 서부, 러시아 동부, 러시아 중부, 그린란드다 [DECISION: 지역 이름 표기] [DECISION: 러시아 중부의 영문 표기]. 알래스카와 주 지역 4곳의 라벨 행 수는 러시아 동부 30개부터 알래스카 13,606개까지다(표~1). 하위 지역 10곳(보충 표~S11)은 라벨 없이 0.5$^{\circ}$ 블록 중심을 $k$-평균으로 군집해 정하였고, 본 실행 전에 고정하였다.

라벨 수 $n$은 행의 수이며, 행의 뜻은 지역마다 다르다. 행은 러시아 서부와 동부에서 CALM 지점의 다년 평균이고, 레나델타에서는 대개 한 번 방문한 관측이며, 알래스카와 캐나다에서는 대부분 점 관측이고, 추가 지역에서는 약 1~km 범위의 평균이다. 1~km 위치당 행 수는 알래스카 39.7, 레나델타 15.1, 캐나다 8.7이며, 알래스카 13,606행은 74개 블록의 343개 위치에 있다.

추가 지역은 어떤 적합보다 먼저 정한 규칙에 따라 공개 자료에서 모았다(보충 방법~2). 라벨은 1990년 이후의 계절 말 융해 깊이를 직접 잰 값이어야 하고, 기존 라벨 지역에서 100~km 이상 떨어져야 하며, 약 1~km 셀로 모았다. 유효 분할의 채점 블록 합집합이 8개 이상인 지역만 신뢰구간 산출에 적격으로 보았다. 재배포 약관이 확인되지 않은 자료는 주 판정을 정하는 약관 확인판에서 뺐다(보충 표~S12). 이 판에서는 러시아 중부(57행, 13블록, 얕은 레짐)와 티베트 고원(132행, 34블록, 깊은 레짐)이 적격이다. 북대서양 묶음(41행)은 전체 판에서만 적격이며 보충 정보에 보고한다.

\subsection{공변량}\label{sec:m-covariates}

모든 학습 방법은 모든 지역에서 쓸 수 있는 공변량 25개를 쓴다. 지형 공변량 6개는 고도, 경사, 사면 방향의 사인과 코사인, $33 \times 33$ 화소 창의 지형 위치 지수와 거칠기이며, Copernicus DEM GLO-30에서 계산하였다(자료 가용성). 기후 공변량 8개는 연평균 기온, 융해 도일(TDD), 동결 도일, $\sqrt{\mathrm{TDD}}$, 최난월과 최한월 기온, 1층 토양 온도, 적설 수량이며, 2015--2020년 ERA5-Land 월평균에서 얻었다\cite{munozsabater2021_era5_land}. 토양 공변량 9개는 5--15~cm 깊이의 점토, 모래, 실트 함량, 용적 밀도, 조립 파편, pH와 세 깊이의 토양 유기탄소이며, 250~m SoilGrids 2.0 층\cite{Poggio2021}을 약 5~km로 재표본해 추출하였다. 모델 기반 공변량 2개는 ESA CCI Permafrost ALT 제품 4.0판의 1997--2021년 평균과 그 유효 표지다\cite{Westermann2024}. 이 제품은 같은 재분석 자료로 일부 구동한 모델 산출물이므로 관측으로 다루지 않는다.

TDD와 Stefan 해\cite{Stefan1891}는 다음과 같다.
\begin{equation}
\mathrm{TDD} = \sum_{m} \max(\bar{T}_m, 0)\, d_m, \qquad \mathrm{ALT} \approx E\, s, \qquad s = \sqrt{\mathrm{TDD}},
\label{eq:stefan}
\end{equation}
여기서 $\bar{T}_m$은 $m$월의 2~m 기온 월평균, $d_m$은 그 달의 일수이고, $E$의 단위는 cm\,($^{\circ}$C\,d)$^{-1/2}$이다. 셀은 라벨, CCI 값, 1층 토양 온도로 계산한 TDD가 모두 유효할 때만 채점하며, 모든 방법이 같은 채점 셀을 쓴다. 관측 횟수나 연도처럼 라벨에서 유도한 양은 입력에서 빼며, 단위 시험으로 이를 강제한다. 입력의 정보 해상도는 기후가 약 0.1$^{\circ}$, 토양이 약 5~km다.

\subsection{Stefan 기준선과 계수 재보정}\label{sec:m-stefan}

계수 $E$는 토양 열 특성, 함수량, 지표 조건을 흡수하며, 지역 적용에서는 현지 측정으로 보정한다\cite{Nelson1997}. 대상의 원천 집합 $S$는 대상 지역과 그 둘레 100~km 완충 구역 밖의 모든 라벨 셀이며, 원천 계수는 원점을 지나는 최소제곱 기울기다.
\begin{equation}
E_0 = \frac{\sum_{i \in S} s_i\, y_i}{\sum_{i \in S} s_i^2},
\label{eq:e0}
\end{equation}
여기서 $y_i$는 관측 ALT다. 원천 계수 Stefan 식은 $E_0 s$를 예측한다. 주 대상의 원천 집합에서 알래스카 셀은 78--94\%를 차지하므로, $E_0$는 이 대상들 사이에서 거의 같다(1.59--1.62~cm\,($^{\circ}$C\,d)$^{-1/2}$; 표~1). 같은 추정량을 대상 라벨 $n$개로 계산한 값을 $E_{\mathrm{ls}}$라 하면, 재보정 Stefan 식은 다음 $E_n$으로 $E_n s$를 예측한다.
\begin{equation}
E_n = \frac{n\, E_{\mathrm{ls}} + \kappa\, E_0}{n + \kappa}, \qquad \kappa = 10.
\label{eq:en}
\end{equation}
따라서 원천 계수는 라벨 10개의 가중을 가지며, 이는 라벨이 적을 때 재보정이 오차를 키울 수 있기 때문이다\cite{Steyerberg2004}. $\kappa$ 값은 이 과제의 앞선 실험에서 정하였으므로 맹검이 아니며, $\kappa = 3$과 $\kappa = 30$은 민감도 사례다. 현지 최소제곱 Stefan 식은 수축 없이 $E_{\mathrm{ls}}$를 쓴다. 기계학습(ML)의 이득은 같은 라벨로 맞춘 재보정 모델에 대해서만 말한다. 원천 계수에 대한 이득에는 재보정의 이득이 들어 있기 때문이다.

물리 모델과 ALT 제품을 비교한 등록 분석에서 더 강한 기준선 두 가지를 더하였다. 연도 정합 Stefan 식은 각 라벨의 관측 연도(2010--2024년 안, 그 이전 라벨은 2010--2014년)에 대해 평균한 $\sqrt{\mathrm{TDD}}$에 $E_0$를 적용하며, 라벨 0개와 전량에서 보고한다. 라벨 40개와 160개에서는 토양 보정 Stefan 식(보충 방법~3)의 출력을 식~(\ref{eq:e0})과 식~(\ref{eq:en})으로 재보정해 더 강한 재보정 기준선으로 보고한다.

\subsection{기계학습 방법과 결합 구조}\label{sec:m-ml}

전이 설계에서 학습 행은 원천 셀과 대상 라벨 $n$개이며 가중은 같다. 앵커를 쓰는 방법은 다음 꼴이다.
\begin{equation}
\hat{y}(x) = a(x) + \lambda\, g(x),
\label{eq:anchor}
\end{equation}
여기서 $x$는 셀의 공변량, $a$는 물리 앵커, $g$는 그 잔차로 학습한 학습기다. 직접 ML은 $\sqrt{\mathrm{TDD}}$와 CCI 값을 포함한 공변량으로 ALT를 회귀한다. 물리 유사라벨 증강은 대상의 라벨 절반에 있는 셀에 값이 $E_n s$인 행을 원천 행 하나당 10개씩 더하며, 추출된 셀은 관측값을 유지한다. 원천 앵커 잔차 결합은 $a = E_0 s$를 쓴다. 재보정 앵커 잔차 결합(잔차 ML)은 $a = E_n s$를 쓰며, 잔차 목표는 원천 행에서 $y - E_0 s$, 대상 행에서 $y - E_n s$다. 증강·앵커·잔차 결합은 라벨 절반의 셀에 값이 0인 잔차 행을 더한다. 물리 입력 ML은 Kudryavtsev형 모델\cite{Kudryavtsev1974}과 토양 보정 Stefan 식의 출력, 또는 앵커 값을 직접 ML의 입력에 더한다. Stefan·CCI 앵커는 CCI 값이 유효한 곳에서 $E_n s$와 그 값의 평균이다. 본 실험의 다른 방법 네 가지는 보충 방법~3에 적었다. 잔차 가중 $\lambda \in \{0.25, 0.5, 1.0\}$은 적합 뒤에 적용한다. 판정의 기준 가중은 본 실행 전에 정하였으며, 잔차 구조에서 0.25, 직접 구조와 물리 입력 구조에서 1.0이다. 가중 0.25의 원천 앵커 잔차 결합을 저가중 잔차라 부른다.

위약 네 가지로 증강의 이득이 셀 수준 Stefan 값에서 오는지 시험하였다. 위약은 유사라벨을 셀 사이에서 섞은 같은 값, 대상 수준 평균, 원천 ALT 평균, 원천에 맞춘 TDD의 선형 함수로 바꾼 것이다.

방법 비교의 학습기는 CatBoost\cite{Prokhorenkova2018}다(반복 200회, 학습률 0.05, 깊이 3, seed 2개). 대상 라벨 없는 직접 ML에서는 학습기 10종을 비교하였다. 이 설정, 더 큰 CatBoost, 원천 지역 하나 제외 교차검증으로 조정한 CatBoost, 랜덤 포레스트, 미세 조정 없는 TabPFN v2\cite{Hollmann2025}와 TabICL v2\cite{Qu2026}, 그리고 원천 자료에서 원천 지역 하나 제외 교차검증으로 초모수를 고른 신경망 4종이다\cite{BergstraBengio2012,CawleyTalbot2010}. 신경망은 다층 퍼셉트론, 다중 헤드 퍼셉트론(TabM\cite{Gorishniy2025}과 다름), 축소 FT-Transformer형 모델\cite{Gorishniy2021}, RealMLP\cite{Holzmuller2024}다(설정은 보충 방법~5). 그림에서는 이 방법들을 Source Stefan, Year-matched Stefan, Recalibrated Stefan, Anchor + residual ML, Source anchor + residual, Physics pseudo-labels, Direct ML, Physics-input ML, Stefan-CCI anchor로, 아래에 적은 규칙은 Target-label CV selection으로 표시한다.

\subsection{전이 시험 설계}\label{sec:m-transfer}

가까운 라벨 셀에서 정보가 새지 않도록, 대상의 원천은 대상 지역과 100~km 완충 구역을 뺀다. 하위 지역은 상위 지역을 원천에서 뺀 방식과, 라벨이 있는 지역 안의 새 구역을 흉내 내도록 상위 지역의 나머지를 남긴 방식으로 평가하였다. 대상 27개는 알래스카를 뺀 3판 지역 6곳, 알래스카, 두 방식의 하위 지역 10곳이다.

각 대상의 0.5$^{\circ}$ 블록을 무작위로 순서를 바꾼 뒤 셀이 적은 쪽 절반에 차례로 배정하여 절반 A와 B로 다섯 번 나누었으므로, 라벨 셀과 채점 셀은 블록 경계로 떨어진다\cite{Roberts2017}. 라벨은 A에서 뽑고 오차는 B의 채점 셀에서 쟀다. A의 크기보다 작은 $\{0, 3, 10, 40, 160, 320, 1000\}$의 각 $n$과 A의 전체 셀에 대해, 대상, 방식, 분할, $n$, 추출로 정한 seed로 다섯 번 무작위 추출하여 모든 방법이 같은 라벨을 쓰게 하였다. 러시아 서부와 동부는 A의 셀이 14--17개이므로, 라벨 40개와 160개의 풀 평균에는 레나델타와 캐나다만 들어간다.

주 지역 4곳(레나델타, 캐나다, 러시아 서부, 러시아 동부)이 확인적 풀이다. 이 지역들은 이 과제의 앞선 실험에 쓰였으므로, 그 판정은 재사용 지역의 재검정이다. 알래스카는 참조 대상으로서 따로 보고하며, 레나델타, 캐나다와 함께 보조 3지역 평균에 넣는다. 3판 표의 러시아 중부와 그린란드는 점 추정으로만 보고하고, 상위 지역과 셀을 공유하는 하위 지역은 판정 범주별로 센다. 추가 지역은 얕은 레짐 5지역 확장 풀과 티베트 고원을 포함한 6지역 확장 풀에 들어간다. 티베트 고원의 채점 셀은 모두 원천 셀 고도의 99.5 백분위수보다 높으므로, 티베트 대비는 독립 지역의 근거로 쓰지 않는다.

검증 사다리는 알래스카와 주 지역 4곳에서 직접 ML과 학습 폴드에 맞춘 Stefan 식을, 무작위 셀부터 지역 홀드아웃까지 설계상 분리가 커지는 여덟 방식으로 채점하였다(보충 방법~3). 앞선 알래스카 분석은 $k$겹 최근접 거리 정합(kNNDM)\cite{Linnenbrink2024}을 더하였는데, 이 방식은 시험과 학습 사이의 거리를 예측 격자와 라벨 사이의 거리에 맞춘다.

\subsection{지역 내 시험 설계와 이득 분해}\label{sec:m-within}

지역 내 시험은 라벨이 수백에서 수천 개인 곳에서 ML이 무엇을 더하는지 묻으며, 전이 결과를 본 뒤 설계하였다. 알래스카, 레나델타, 캐나다에서 라벨은 절반 A에서 뽑았고 유일한 학습 행이었다. 전이 설계의 원천 계수를 사전값으로 썼고, $\lambda$는 뽑은 라벨 안의 블록 교차검증으로 $\{0.25, 0.5, 1.0\}$에서 골랐으며, 증강은 라벨 하나당 유사라벨 10개를 썼다. 첫 시험(분할 5개, 라벨 1000개 이상, 비교 상대는 교차검증으로 고른 최선의 물리 보정식, 알래스카에서는 SAR 공변량 9개 추가)은 기각되었다. 두 번째 시험은 네 요소를 바꾸었다. 비교 상대(재보정 Stefan 식), 분할(25개, 레나델타는 24개 유효), 레시피(Stefan·CCI 앵커와 직접 ML 추가), 라벨 범위(200, 500, 1000개와 전량)다. 두 시험을 모두 보고한다.

이득 분해는 ML이 어느 공간 규모에서 오차를 줄이는지 묻는다. 채점 셀을 0.5$^{\circ}$ 블록과 기온 $\sqrt{\mathrm{TDD}}$로 묶어, 한 묶음 안에서 Stefan 예측과 기후 공변량 8개가 일정하게 하였다. 묶음 $g$의 셀 수를 $N_g$, 평균을 $\bar{y}_g$와 $\bar{\hat{y}}_g$라 하면 다음과 같다.
\begin{equation}
\mathrm{SSE} = \sum_g \sum_{i \in g} \left[(y_i - \bar{y}_g) - (\hat{y}_i - \bar{\hat{y}}_g)\right]^2 + \sum_g N_g \left(\bar{y}_g - \bar{\hat{y}}_g\right)^2,
\label{eq:decomp}
\end{equation}
여기서 두 항은 격자 안 성분과 격자 사이 성분이다. 격자 안 RMSE는 셀이 2개 이상인 묶음을 쓰며, Stefan 식에서는 관측의 묶음 내 표준편차와 같다. 성분의 RMSE 차이는 더해지지 않으므로, 감소의 비율은 SSE로 계산한다. 등록 문서는 묶음 기준으로 토양 $\sqrt{\mathrm{TDD}}$를 적었으나 실행은 Stefan 예측의 기온 $\sqrt{\mathrm{TDD}}$를 썼다. 그래서 토양 기준 묶음의 민감도 분석을 등록하고 실행하였다(두 묶음 사이 수정 Rand 지수는 알래스카 0.999, 캐나다 0.997, 레나델타 0.832).

\subsection{라벨 배치 절차}\label{sec:m-placement}

배치 전략은 능동 선정을 빼고는 라벨 값을 쓰지 않고 절반 A의 셀을 고른다. 셀 무작위 추출 외의 전략은 블록 층화, 공변량 군집 층화, 공변량 최대 최소 거리 추출, 원천 외삽 우선, 잔차 학습기의 예측 분산에 따른 능동 선정, $\sqrt{\mathrm{TDD}}$ 우선이다(보충 방법~8). 지역 안에서는 라벨 20--500개에서 비교하였고, 블록 층화와 최대 최소 거리 추출은 전이에서 라벨 10개와 40개로도 비교하였다. 이 시험들은 결과를 본 뒤 설계하였고, 본 실험의 블록 분산 대비는 결과 전에 등록하였다.

알고리즘 P는 블록 배분과 공변량 범위를 결합한 절차이며, 워크플로 분석에서 시험한 사전 지정 배치 절차다. 입력은 후보 셀의 표준화 공변량(라벨 값 없음), 블록 크기 $N_b$, 라벨 10, 40, 160개의 중첩 예산, 배분 지수 $\gamma$, seed다.
\begin{enumerate}
\item seed로 고른 블록이나 기존 라벨이 있는 블록에서 시작해, 블록의 공변량 중심 $\mu_b$를 가장 먼 점 우선 순회로 정렬한다.
\item 예산 $n_k$에서 블록 $b$에 $q_b = n_k N_b^{\gamma} / \sum_{b'} N_{b'}^{\gamma}$개의 라벨을 최대 나머지 방식으로 배분하되 $N_b$를 넘지 않게 한다.
\item 남은 부족분이 가장 큰 블록을 채운다. 그 블록에 라벨이 없으면 $\mu_b$에 가장 가까운 셀을, 있으면 이미 고른 모든 라벨에서 가장 먼 셀을 고른다.
\item 예산이 늘어도 앞서 고른 라벨은 유지한다.
\end{enumerate}
후보 설정은 $\gamma = 0$, 0.5, 1, 블록의 첫 셀을 가장 먼 점 우선으로 고르는 변형, 블록 층화, 최대 최소 거리 추출이다. 평가 지역이 선택에 영향을 주지 않도록 알래스카 대상만으로 기계적 규칙에 따라 후보 하나를 고른다. 이 규칙에서 모든 후보가 제외되어 Algorithm P는 셀 무작위 추출이며, Algorithm P 대 무작위 배치의 워크플로 대비는 시험하지 않았다. 따라서 워크플로의 배치 단계는 기여한 것이 없다. 라벨 10개에서는 워크플로와 그 무작위 배치 비교 대상이 같은 라벨과 모델을 쓰고, 라벨 40개와 160개에서 무작위 배치와 고정 잔차 레시피에 대한 비교는 같은 라벨에서 선정 규칙과 고정 레시피를 비교한 것이다. 등록한 보고 규칙에 따라 측정 위치의 순위를 매긴 지도는 만들지 않았다.

\subsection{방법 선정 규칙과 편향 진단}\label{sec:m-selection}

선정 규칙은 대상 라벨만으로 방법을 고른다.
\begin{enumerate}
\item 라벨이 10개보다 적으면 재보정 Stefan 식을 쓴다.
\item 라벨이 있는 블록을 최대 5개 폴드에 무작위로 배정한다(블록이 하나면 셀 폴드). 폴드가 2개보다 적으면 재보정 Stefan 식을 쓴다.
\item 각 폴드에서 $E_n$을 다시 계산하고 후보 7개를 맞춘다. 후보는 원천 계수, 재보정, 현지 최소제곱 Stefan 식, $\lambda = 0.25$와 1.0의 잔차 ML, 증강·앵커·잔차 결합, 물리 유사라벨이다.
\item 교차검증 RMSE가 가장 낮은 후보를 고른다(동률이면 작은 $\lambda$와 앞선 후보).
\item 고른 후보를 모든 라벨로 다시 맞춘다.
\end{enumerate}
이 규칙은 라벨 격자 결과를 본 뒤 등록하였고, 대상 30개(전이 대상 27개와 약관 확인 추가 지역 대상 3개)의 라벨 10, 40, 160개와 전량에 적용하였다. 비교 상대는 고정 잔차 레시피이며 비열등 한계는 0.5~cm다.

편향 진단값은 한 추출의 라벨 10개에서 원천 계수 Stefan 식이 보인 평균 오차의 절댓값을 추출과 분할에 걸쳐 평균한 값이다. 재보정도 쓰는 라벨만 쓰므로, 어떤 학습기를 맞추기 전에 알 수 있다. 이 진단값과 라벨 전량에서 선정 규칙이 원천 계수 모델에 대해 보인 이득의 Spearman 상관을 대상 28개에서 계산하였다. 이 이득에는 재보정이 들어 있으므로, 상관은 라벨을 쓰는 보정(재보정과 잔차)에 관한 것이다. 재보정을 넘어선 이득의 몫은 따로 사후 분석으로 살폈다(추가 등록 분석).

\subsection{통계 분석}\label{sec:m-stats}

절반 B의 블록 $b$에서 오차 제곱합을 $\mathrm{SSE}_b$, 채점 셀 수를 $N_b$라 할 때, 두 가지 평균 제곱근 오차(RMSE, cm)를 계산하였다.
\begin{equation}
\mathrm{RMSE}_{\mathrm{cell}} = \left(\frac{\sum_b \mathrm{SSE}_b}{\sum_b N_b}\right)^{1/2}, \qquad \mathrm{RMSE}_{\mathrm{block}} = \frac{1}{|B|}\sum_b \left(\frac{\mathrm{SSE}_b}{N_b}\right)^{1/2},
\label{eq:rmse}
\end{equation}
여기서 $|B|$는 채점 블록 수다. 셀 가중 값이 주 값이고, 블록 등가중 값은 표본이 조밀한 블록의 영향을 줄인다. 대비 $\Delta$는 첫 방법의 RMSE에서 둘째 방법의 RMSE를 뺀 값으로, 분할, 추출, seed로 짝지어 분할에 걸쳐 평균하므로 음수는 첫 방법의 오차가 낮다는 뜻이다. 층화 평균은 지역에 같은 가중을 준다.

신뢰구간(CI)은 각 분할 안에서 채점 블록을 부트스트랩한 95\% 백분위 구간이며, 두 방법에 같은 재표집 블록을 쓴다. 라벨 격자 하네스가 모은 가설(본 실험과 추가 지역 실행)과 예측 구간 실험은 반복 1000회를, 다른 모든 대비는 10,000회를 썼다. 두 팔이 서로 다른 라벨 집합을 쓰는 본 실험의 블록 분산 대비와 워크플로 분석은 추출도 재표집하는 2단 부트스트랩을 썼고, 배치 시험은 라벨 추출을 조건으로 한 구간을 보고한다. 보조 구간은 한 지역 블록의 합집합을 모든 분할에 대해 한 번 재표집하였다. 이 구간이 판정을 바꾸는 경우 결과를 분할 독립 가정에 의존한다고 적으며, 이 단서가 없는 문장은 두 판정이 일치해야 한다. 지역은 유효 분할의 채점 블록 합집합이 8개 이상일 때 풀 구간에 들어갔다.

대비는 셀 가중과 블록 등가중 CI의 상한이 모두 0보다 작으면 오차 감소, 하한이 모두 0보다 크면 오차 증가, 둘 다 아니고 네 한계가 모두 $\pm 0.5$~cm 안에 있으면 동등(유의수준 2.5\%의 두 단측 검정\cite{Schuirmann1987}), 그 밖에는 미결정('차이를 확인하지 못했다')이다. 한계는 어떤 결과보다 먼저 정하였다. 이 한계는 주 지역에서 원천 계수 Stefan 식의 셀 가중 RMSE(21.7--43.0~cm)의 약 1--2\%이며, 측정 오차가 아니라 관례다. 1.0~cm와 그 RMSE의 2\%를 한계로 한 경우는 민감도 사례다. $|\Delta| < 0.5$~cm인 오차 감소나 오차 증가는 크기가 0.5~cm 미만인 구별 가능한 차이로 보고한다. 비열등은 단측이며 두 상한이 모두 $+0.5$~cm보다 작아야 한다. 대비의 최소 라벨 수 $n^{*}$는 셀 가중 CI 상한이 0보다 작고, 분할의 3분의 2 이상과 추출의 75\% 이상에서 $\Delta < 0$인 가장 작은 $n$이다.

유의수준은 양측 $\alpha = 0.05$이며, 단측 0.025에 해당한다. $P$ 값은 두 가중의 양측 부트스트랩 $P$ 값 가운데 큰 값이며, 모든 등록 대비의 보정 전과 Holm 보정 $P$ 값은 보충 표~S1에 있다. 결과 전에 정한 대비 10개만 초록에서 판정어를 쓸 수 있다. 이들의 $P$ 값은 가족 크기를 10으로 고정한 Holm 절차\cite{Holm1979}로 보정하였고, 보정 $P$가 0.05보다 작을 때만 초록에 방향을 적는다. 그렇지 않으면 초록에는 차이를 확인하지 못했다고 적고, 본문에 '보정 전에만 유의'를 더한다. 지역 내 시험, 배치 시험, 선정 시험과 추가 분석의 결과는 초록에 수치로만 나온다. 다른 등록 가족에서 Holm 보정 $P$ 값은 보조 값이며, 워크플로 시험과 제품 비교에서만 문구를 정한다.

풀 대비는 4--7개 지역에 근거하므로 재표집 단위는 0.5$^{\circ}$ 블록(지역당 13--74개)이며, 모든 풀 평균에 지역 값과 범위를 함께 적는다. 지역 일반에 관한 문장을 위해 핵심 대비를 주 지역 4곳과 알래스카의 분할 50개로 다시 계산하고, 부호 검정, 부호 뒤집기 순열 검정, Hartung-Knapp 확률효과 평균\cite{HartungKnapp2001}으로 요약하였다. 이런 문장은 Hartung-Knapp 구간이 0을 벗어날 때만 쓰며, 지역이 4곳이면 가장 작은 양측 부호 검정 $P$는 0.125다. 순위상관은 Spearman 계수이며, 구간은 블록과 추출의 결합 재표집으로 얻었다. 오차 하한은 지형이 아닌 거친 공변량이 같은 셀 5개 이상의 묶음 안에서 관측의 합동 분산이며(알래스카 13,333셀), 효과 크기의 분모로만 쓴다.

앞선 결과를 본 뒤 설계한 분석은 그렇게 표시하고, 사후 서술에는 판정어를 쓰지 않는다. 등록 가설마다 맹검 표지 네 가지 가운데 하나를 단다. 맹검, 부분 비맹검(다른 규약에서 방향을 봄), 재현(비맹검)(같은 대비의 방향을 봄), 등록 시점에 결과 존재(열지 않음)다. 예측 구간은 경험적 포함률과, $a = 0.1$인 중앙 구간 $[L, U]$의 구간 점수 $(U - L) + (2/a)(L - y)_{+} + (2/a)(y - U)_{+}$로 평가하였다\cite{GneitingRaftery2007}. 계층 conformal 구간은 원천 지역을 묶음으로 썼으며\cite{Dunn2023}, 묶음이 2--4개이므로 유한 표본 보장이 없다(보충 방법~4). $\pm$ 뒤에 값을 적지 않는다.

\subsection{내부 사전 등록과 등록 이탈}\label{sec:m-registration}

가설, 대비, 판정 규칙, 허용 문장은 해당 결과를 가져오기 전에 과제의 git 저장소에 커밋하였다(내부 사전 등록). 커밋에는 제삼자 시간 인증이 없으며, 시각은 한국 표준시의 커밋 시각이다(전체 표는 보충 방법~6). 라벨 격자 실험(LG, L1--L8)은 2026년 9월 29일 15:45:46에 커밋 40be64c로 등록하였다. 그 확장(LGX), 파운데이션 모델과 추가 지역 실행(LGT, LGD), 구간 실험(LGU), TabICL과 조정 신경망 실행(LGF), 초록 묶음 AB1--AB10을 포함한 보고 규칙(WRAPUP, 316714c, 2026년 9월 30일 03:02:42)이 뒤따랐고, 모두 그날 07:20:22의 첫 결과 회수보다 앞섰다.

워크플로 실험은 이 결과를 본 뒤 설계하고 실행 전에 등록하였다. WF0--WF5는 4168331(2026년 10월 1일), WF6--WF8은 1b42b4d(2026년 10월 1일), WF9와 WF10은 8180632(2026년 10월 2일), 토양 기준 묶음 민감도는 7da0064(2026년 10월 2일)다. WF0은 사후 서술이고, WF5로 계획한 우선순위 지도는 만들지 않았다. 분석 XA--XJ는 어떤 적합보다 먼저 커밋 2678100(2026년 10월 4일 14:22:39)으로 등록하였다. 실행 전 등록 확인 시험으로 표시한 것은 XC-F3뿐이며, XA는 사후 분석이다. 각 분석을 본문에 둘지 보충 정보에 둘지는 결과 전에 정하였다.

등록 이탈과 사후 변경은 다음과 같다.
\begin{enumerate}
\item 이득 분해(WF9)는 등록한 토양 $\sqrt{\mathrm{TDD}}$가 아니라 기온 $\sqrt{\mathrm{TDD}}$로 셀을 묶었고, 토양 기준 묶음의 민감도 분석을 더하였다.
\item 첫 지역 내 시험(WF1-a)이 기각된 뒤 두 번째 시험(WF6)은 비교 상대, 분할, 레시피, 라벨 범위를 바꾸었다.
\item 제품 비교(XG)는 구간 가설(LGU-B2)에 조건을 둔 서술 표로 등록되었다. 그 가설이 미결정으로 끝난 뒤 XG는 Holm 보정 대비 6개로 다시 등록하였고 등록 이탈 표지를 단다.
\item 워크플로 분할(XC)은 seed 1--200이 이미 쓰였으므로(LGX-N4) seed 6--15에서 201--210으로 옮겼고, XC 판정어를 초록에 쓰자는 제안은 철회하였다.
\item AB4에는 결과 전 민감도 계산이 같은 양의 분할 분포를 보여 주었다는 주석을 단다.
\item 표를 고정한 뒤 발견한 지온 유도 셀 20개는 표지를 단 채 직접 라벨 범주에 남아 있다.
\item 기술적 수정 네 건(ff1be02, 2026년 9월 30일 11:44:30)은 첫 회수 뒤에, 그러나 어떤 결과도 풀거나 열기 전에 커밋하였으며, 가설, 대비, 판정 규칙을 바꾸지 않는다. [DECISION: 수용 여부, 방법 초안 183행]
\item CatBoost는 교차 환경 점검 (i)을 통과하지 못했으므로, 플랫폼을 섞은 확장 풀 행에는 교차 환경 표지를 단다.
\item 등록 규칙에 따라 측정 우선순위 지도는 만들지 않았다.
\item 기후 외삽 시험(WF10)은 캐나다의 외삽 집합이 블록 3개여서 판정할 수 없었다.
\item XA는 등록한 7개 계열이 아니라 8개 대상 계열로 재표집하였다. 28대상 집합의 북대서양 대상이 등록 계열 목록에 없었기 때문이다(봉인 결과와 함께 기록한 구현 이탈).
\end{enumerate}

\subsection{추가 등록 분석}\label{sec:m-xbatch}

분석 10개(XA--XJ)는 위의 규약을 따르며, 각 분석을 본문과 보충 정보 가운데 어디에 둘지는 결과 전에 정하였다(보충 방법~9; 결과는 보충 표~S2).

XA(본문, 사후)는 원천 계수 모델에 대한 이득을 재보정 몫, 남은 수축 몫, 재보정을 넘어선 ML 몫으로 나누고, 두 정의의 계수 오차와 각 몫의 순위상관을 추정하며, 구간은 대상 계열과 블록의 결합 재표집으로 얻는다. 등록 목록은 7개 계열이었고, 그 목록에 없는 북대서양 대상은 8번째 단일 대상 계열로 재표집하였으며, 등록한 7개 계열의 27대상은 결과 범주가 같은 등록 밖 민감도로 분석하였다. 결과는 0을 벗어난 구간이거나 '확인하지 못함'으로만 적는다. XB(보충)는 교차 적합으로 대상 라벨에서 물리 모델과 ALT 제품의 볼록 가중을 학습하며, 잔차 학습기가 있는 경우와 없는 경우를 본다. XC(본문)는 알고리즘 P, 라벨 10개 진단, 방법 규칙을 얕은 레짐 5지역의 새 분할에 차례로 적용한다. Holm 보정 가설 8개는 라벨 40개와 160개에서 워크플로를 무작위 배치와 고정 잔차 레시피에 비교하고, 라벨 10, 40, 160개에서 재보정 Stefan 식에 대한 비열등을 시험하며, 선정 규칙 아래에서 알고리즘 P를 무작위 배치와 비교하고, 새로 얻은 지역에서 라벨 10개의 워크플로를 시험한다. 평가 셀이 구성 요소 설계에 쓰였으므로, 마지막 가설을 뺀 모두가 부분 비맹검이다.

XD(보충)는 알래스카 대상에서 알고리즘 P의 설정을 고르고, 다른 계수 추정량을 시험하며, 학습된 배치 정책을 알래스카 밖에서 $P$ 값 없이 한 번 평가한다. XE(보충)는 지역 안 잔차 ML에 10--500~m의 공개 입력을 더하고, 격자 안 RMSE를 공변량 25개 모델과 재보정 Stefan 식에 비교한다. XF(보충, 맹검)는 새 공개 자료에 적격 규칙을 적용하고 적격 지역마다 핵심 대비를 되풀이한다. XG(본문)는 제품 학습 지점에서 5~km나 25~km 안의 블록을 가린 뒤 기존 ALT 지도를 같은 블록에서 채점하며, 원래 값은 원천 계수 모델에, 재보정 값은 잔차 ML에 비교한다. XH(본문)는 세 지역에서 무작위 셀, 지점, 블록, kNNDM, 지역 홀드아웃 검증에 따른 재보정 Stefan 식, 잔차 ML, 직접 ML의 오차를 판정어 없이 서술한다. XI는 약관 확인 캐나다 확충 표에서 기후 외삽 시험을 되풀이하고, XJ는 지온 유도 라벨을 저가중 보조 행으로 시험한다(둘 다 보충).

\subsection{계산 환경과 재현성}\label{sec:m-computing}

물리 모델과 CatBoost는 CPU에서 결정적으로 실행된다. 라벨 격자 실험과 그 확장은 CPU 64코어와 NVIDIA T4 GPU 4개를 가진 클라우드 노드(Rescale)에서, 워크플로 실험과 추가 분석은 한 프로세서 계열(AMD EPYC Milan 96코어, Milan-X 64코어)의 클라우드 CPU 노드에서 실행하였다. 파운데이션 모델, 조정 신경망, 구간 실험, 추가 지역 실행은 NVIDIA RTX 3090 GPU를 가진 공유 로컬 서버를 썼다. 패키지 판은 고정하였고(모든 플랫폼에서 CatBoost 1.2.10; 보충 표~S13), 각 산출물은 코드 해시, 설정 해시, 패키지 판, 장치를 기록한다.

물리식과 CatBoost 결과는 클라우드와 로컬 실행 사이에서 34,360개 키에 걸쳐 $3.6 \times 10^{-14}$~cm 안에서 일치하였고, 클라우드 노드 두 종류는 공통 키 11,814개에서 같은 오차를 냈다. 신경망은 플랫폼 사이에서 재현되지 않았다. FT-Transformer 키 288개 가운데 43개가 0.5~cm 넘게 달랐고(최대 5.61~cm), 한 플랫폼 안의 재적합은 같았다. 릿지 앵커 위의 CatBoost는 최대 0.46~cm 달랐다. 따라서 대비의 양쪽은 한 플랫폼에서 오고, 플랫폼을 넘는 풀 평균에는 교차 환경 표지를 달며, 워크플로는 CatBoost와 닫힌 형태의 물리식만 쓴다.

\subsection{지도와 색표}\label{sec:m-maps}

지도는 70$^{\circ}$~N에서 축척이 참인 북극 평사 투영을 쓰고 중심 경도는 각 설명문에 적으며, 티베트 고원은 Lambert 방위 정적 투영의 삽도로 넣었다. 육지는 Natural Earth(1:50m), 영구동토 구역은 ESA CCI Permafrost 범위 제품 4.0판(영구동토 비율, 1997--2021년 평균)에서 가져왔고, 지도는 Cartopy 0.23.0으로 그렸다. 과학 색표\cite{Crameri2018Zenodo,Crameri2020}를 썼다. ALT에는 oslo를 뒤집어 0.12--0.90 범위로, 오차 차이에는 0을 중심으로 한 broc을, 잔차에는 bam을, 구간 폭에는 acton을 뒤집어 0.05--0.95 범위로 썼다.

\subsection{대형 언어 모델 사용}\label{sec:m-llm}

[DECISION: LLM 사용 기재 문구]

\subsection{코드 가용성}\label{sec:m-code}

자료 준비, 실험, 통계, 그림, 시험에 쓴 코드는 [DECISION: 공개 저장소 주소와 공개 시점]에서 받을 수 있고, 판 DOI와 함께 보관한다 [DECISION: Zenodo DOI]. 주요 구성 요소는 라벨 격자 실험(\path{scripts/3_deep_learning/h40_label_grid.py}), 검증 사다리(\path{scripts/2_evaluation/h41_validation_ladder.py}), 확장 실험(\path{scripts/3_deep_learning/h42_label_grid_ext.py}), 워크플로 실험(\path{scripts/3_deep_learning/h54_workflow.py}), 초록 대비 집계기(\path{scripts/2_evaluation/h39_scenarios.py}), 사후 재분석(\path{scripts/2_evaluation/wf0_reanalysis.py}), 추가 분석(\path{x_*.py}), 단위 시험(\path{tests/}), 클라우드 작업 설정(\path{configs/rescale/})이며, 등록 문서와 그 개정 이력도 함께 둔다. 추가 분석은 catboost 1.2.10, scikit-learn 1.9.1, pandas 2.3.3, scipy 1.17.1, numpy 2.1.1을 고정한다. TabPFN v2와 TabICL v2의 사전 학습 가중치는 재배포하지 않으며 제공처에서 받는다.
